\documentclass[11pt,a4paper]{article}
\usepackage[margin=2.2cm]{geometry}
\usepackage[T1]{fontenc}
\usepackage{lmodern}
\usepackage{xcolor}
\usepackage{booktabs}
\usepackage{longtable}
\usepackage{enumitem}
\usepackage{hyperref}
\usepackage{microtype}
\hypersetup{colorlinks=true,linkcolor=blue!55!black,urlcolor=blue!65!black}
\setlist{nosep,leftmargin=*}

\title{Pilot Fabric 0.38\\Universal Reliable Television Operation}
\author{Tessaris}
\date{30 August 2026}

\begin{document}
\maketitle

\section{Objective}

Pilot Fabric 0.38 completes the reusable transaction and qualification framework around restricted-device
control.  The transaction implements the sequence

\begin{center}
\textbf{observe $\rightarrow$ act $\rightarrow$ observe $\rightarrow$ verify}
\end{center}

instead of equating command delivery with success.  The first live integration
wraps private-phone navigation, media, pointer and application operations on the
paired LG gateway.  The same contract is device-neutral and can be applied to
Samsung, Android/Google TV, set-top boxes and installed AION nodes as their
adapters become available.

\section{Transaction contract}

Each \nolinkurl{pilot.verified-navigation.transaction.v1} record binds:

\begin{itemize}
  \item the device node and current surface belief;
  \item a normalized user goal and explicit success criterion;
  \item an ordered set of bounded routes and the selected route index;
  \item pre-action device belief and optional image evidence hash;
  \item each action receipt, attempt number and transport result;
  \item structured post-action evidence; and
  \item final state, completion time and canonical transaction hash.
\end{itemize}

Accepted verification criteria are expected surface, forbidden view, changed
screen, OCR text containment, a receipt with governed read-after-write evidence,
and an exact/present device-state field.  This
keeps success logic outside the language model.

\section{State machine}

\begin{longtable}{p{0.27\linewidth}p{0.65\linewidth}}
\toprule
\textbf{State} & \textbf{Meaning} \\
\midrule
\endhead
\nolinkurl{ready_to_act} & A route is selected but no device action receipt has
been accepted. \\
\nolinkurl{awaiting_after_observation} & Transport was delivered, but the
visible or semantic result still needs evidence. \\
\nolinkurl{recovery_ready} & The previous result failed and a predeclared
alternative route remains within the attempt limit. \\
\nolinkurl{verified} & The explicit criterion matched permitted structured
device state or an after-observation. \\
\nolinkurl{not_verified} & Evidence did not establish the goal and no bounded
route remains. \\
\nolinkurl{superseded_not_verified} & A newer user operation replaced an
unfinished operation; the earlier transaction remains evidence rather than
disappearing. \\
\bottomrule
\end{longtable}

The engine accepts at most three device attempts.  It never invents a route
after failure, loops indefinitely, or permits a language model to redefine the
success criterion during execution.

\section{Evidence hierarchy}

Device-state criteria are preferred when the gateway exposes a relevant field.
For example, launching Netflix can be verified against the foreground
application identifier without capturing protected programme pixels.  A remote
directional or profile-selection operation is different: the pointer/button
receipt proves delivery only, so the transaction waits for a structured phone
observation.

Manual capture and explicit Observer Mode process images locally.  The verifier
receives surface, view, observation identifier and hashes, but the source frame
is deleted.  Profile selection normally verifies when the surface remains Netflix
and the profile-chooser view is absent. If OCR cannot identify the changed surface,
an explicit owner visual confirmation may close the transaction only when Pilot
already holds a verified transport receipt and a distinct after-image hash.
Consent dismissal verifies when the consent
view is absent.  Generic focus movement requires distinct before and after image
hashes; a missing before hash cannot be converted into a success.

An in-application Netflix search is verified only when OCR from the subsequent
owner-authorized observation contains the requested title.  Application launch
uses the foreground application identifier.  Play, pause and other third-party
operations require a changed structured screen or signed provider playback
telemetry; an LG transport acknowledgement alone is excluded.

\section{Idempotency and duplicate suppression}

Every private-controller request carries a browser-generated request identifier.
The mother checks the active transaction, last transaction and bounded history
before dispatch.  A repeated identifier returns the original transaction with
\nolinkurl{duplicate_suppressed=true}; it cannot emit a second television action.
This local proof is independently eligible for the duplicate-suppression
qualification area and contains no provider credential or private screen frame.

\section{Recovery and route learning}

Routes are predeclared bounded commands.  When visual verification fails, the
engine may expose the next route as \texttt{recovery\_ready}.  The phone shows a
\emph{Try verified recovery route} button.  Executing it remains inside the same
transaction and again requires post-action evidence.  This prevents a retry from
silently resetting the goal, history or attempt count.

The mother brain stores route statistics under a canonical key derived from
device, surface, goal and route identifier.  Successful and failed counts,
last result and update time are retained.  On a later transaction, routes with
established successes rank ahead of routes with failures.  The memory never
contains a television password, provider credential or source camera frame.

\section{Private-phone interface}

The controller displays the active or most recent goal, transaction state,
attempt count, expected criterion and the number of remembered routes.  When an
after-observation is required it instructs the user to start the explicit
Observer Mode or take a picture.  A recovery button appears only when the engine
has a bounded alternative.  This makes a previously hidden technical distinction
visible in simple language: command delivered does not mean screen verified.
The controller includes governed Back and Home keys, a Netflix title field,
named household profiles and a qualification panel.  The user-facing control is
labelled emph{Switch profile}. Pilot learns the visible profile labels from an
owner-authorized chooser observation, maps each label to its observed position,
and retains no Netflix credential. The field flow is:

\begin{center}
\textbf{capture before $\rightarrow$ act once $\rightarrow$ capture after
$\rightarrow$ inspect the terminal result}
\end{center}

The same procedure qualifies directional movement, focus/pointer actions,
Back/Home/Exit and playback without allowing the interface to self-certify.

\section{Field qualification contract}

Section 3 is closed per television adapter, not globally.  Closure requires at
least twenty eligible field cases, a verified-or-bounded-recovery rate of at
least 95\%, and successful evidence in all eleven areas: connection recovery,
state observation, reversible audio, application launch, directional navigation,
Back/Home/Exit, focus/pointer, in-application search, profile selection, playback
control and duplicate suppression.

Eligible evidence is limited to device read-after-write, structured local screen
observation, bounded owner visual confirmation, signed provider telemetry and
local idempotency proof. Owner confirmation additionally requires machine evidence
of transport delivery and a changed observer hash. Simulated
contracts and transport-only receipts are retained for engineering diagnostics
but excluded from the denominator and from coverage.  The paired LG adapter
completed 24 eligible field cases with 24 verified outcomes, a 100\% success-or-
recovery rate and successful coverage of all eleven required areas. Section 3 is
therefore closed for this adapter.

\section{Auditability and safety}

Action audits record transaction identifier, television node, goal, state and
transaction hash.  Observation audits record the resulting state and that the
source frame was not retained.  The existing signed capability capsule remains
mandatory for every underlying operation.  Navigation does not broaden power,
purchase, account, legal-consent or communication authority.

\section{Verification}

All 291 Fabric tests pass in release 0.38.  Automated tests establish that
transport delivery waits for observation, a
semantic after-state can verify the goal, failure selects a declared alternative,
unchanged or missing image evidence cannot prove a screen change, device fields
can close pixel-free operations, OCR must contain the requested search title,
duplicate request identifiers suppress redispatch, qualification excludes
simulation and transport-only evidence, owner confirmation cannot bypass missing
machine evidence, profile labels are extracted only from the chooser region, and
learned successful routes are ranked first.
The complete Fabric suite additionally covers webOS controls, the phone security
boundary, observer tokens, local perception, audit integrity and all earlier
capability layers.

\section{Platform scope and remaining field work}

The implemented adapter and current field evidence apply to the paired LG webOS
television.  Samsung, Sony/Google TV and other vendors require their own signed
capability adapters and must pass the same qualification contract.  Unit tests
and protocol simulators can validate parsing, state machines, replay protection
and failure handling, but cannot substitute for eventual real-device evidence.
The LG result closes only this adapter. The next platform work is to implement and
qualify other vendor adapters and to expand bounded recovery routes only when new
field failures provide concrete evidence for them.

\end{document}
